% Options for packages loaded elsewhere
\PassOptionsToPackage{unicode}{hyperref}
\PassOptionsToPackage{hyphens}{url}
\PassOptionsToPackage{dvipsnames,svgnames,x11names}{xcolor}
\documentclass[
  10pt,
  fontset=fandol]{ctexart}
\usepackage{xcolor}
\usepackage[margin=16mm]{geometry}
\usepackage{amsmath,amssymb}
\setcounter{secnumdepth}{-\maxdimen} % remove section numbering
\usepackage{iftex}
\ifPDFTeX
  \usepackage[T1]{fontenc}
  \usepackage[utf8]{inputenc}
  \usepackage{textcomp} % provide euro and other symbols
\else % if luatex or xetex
  \usepackage{unicode-math} % this also loads fontspec
  \defaultfontfeatures{Scale=MatchLowercase}
  \defaultfontfeatures[\rmfamily]{Ligatures=TeX,Scale=1}
\fi
\usepackage{lmodern}
\ifPDFTeX\else
  % xetex/luatex font selection
\fi
% Use upquote if available, for straight quotes in verbatim environments
\IfFileExists{upquote.sty}{\usepackage{upquote}}{}
\IfFileExists{microtype.sty}{% use microtype if available
  \usepackage[]{microtype}
  \UseMicrotypeSet[protrusion]{basicmath} % disable protrusion for tt fonts
}{}
\makeatletter
\@ifundefined{KOMAClassName}{% if non-KOMA class
  \IfFileExists{parskip.sty}{%
    \usepackage{parskip}
  }{% else
    \setlength{\parindent}{0pt}
    \setlength{\parskip}{6pt plus 2pt minus 1pt}}
}{% if KOMA class
  \KOMAoptions{parskip=half}}
\makeatother
\setlength{\emergencystretch}{3em} % prevent overfull lines
\providecommand{\tightlist}{%
  \setlength{\itemsep}{0pt}\setlength{\parskip}{0pt}}
\usepackage{amsmath,amssymb,mathtools,needspace}
\xeCJKDeclareCharClass{CJK}{"2032,"0400->"04FF,"2160->"216B,"2460->"2473,"25A0,"25C7,"25CB,"25CF}
\IfFontExistsTF{Arial Unicode MS}{\xeCJKsetup{AutoFallBack=true}\setCJKfallbackfamilyfont{\CJKrmdefault}{Arial Unicode MS}}{}
\setcounter{secnumdepth}{0}
\setlength{\emergencystretch}{2em}
\allowdisplaybreaks
\usepackage{bookmark}
\IfFileExists{xurl.sty}{\usepackage{xurl}}{} % add URL line breaks if available
\urlstyle{same}
\hypersetup{
  pdftitle={差分方法的收敛性},
  colorlinks=true,
  linkcolor={Maroon},
  filecolor={Maroon},
  citecolor={Blue},
  urlcolor={Blue},
  pdfcreator={LaTeX via pandoc}}

\title{差分方法的收敛性}
\author{}
\date{}

\begin{document}
\maketitle

\subsubsection{5.7.4 差分方法的收敛性}\label{ux5deeux5206ux65b9ux6cd5ux7684ux6536ux655bux6027}

现在运用极值原理论证差分方法的收敛性并估计误差。

\textbf{定理 5.4} 设 \(y_n\) 是差分问题式（5.7.10）的解，而 \(y(x_n)\) 是边值问题（5.7.9）的解 \(y(x)\) 在节点 \(x_n\) 的值，则截断误差 \(e_n=y(x_n)-y_n\) 有下列估计式：

\[
|e_n|\le\frac{M(b-a)^2}{96}h^2,\qquad
M=\max_{a\le x\le b}|y^{(4)}(x)|.
\tag{5.7.11}
\]

\textbf{证明} 由 Taylor 展开，易得

\[
\begin{cases}
\displaystyle\frac{y(x_{n+1})-2y(x_n)+y(x_{n-1})}{h^2}-q_ny(x_n)=r_n+\frac{h^2}{12}y^{(4)}(\xi_n),\quad x_{n-1}<\xi_n<x_{n+1},\\[6pt]
y(x_0)=\alpha,\quad y(x_N)=\beta.
\end{cases}
\tag{5.7.12}
\]

\begin{quote}
校注（agent 补充）：定理 5.2 的负最小值一侧，原书确实写作严格条件 \(l(y_n)<0\)；定理 5.3 却在 \(l(y_n)=0\) 时同时引用正最大值、负最小值结论。此处原文均照录。补足该引用的方法是对 \(-y_n\) 应用定理 5.2 的 \(l\ge0\) 一侧（利用 \(l\) 的线性），从而排除负最小值；无需改动原书方程。
\end{quote}

\href{../../source-images/pdf-152.jpeg}{原页}

将式（5.7.12）与式（5.7.10）相减，知截断误差 \(e_n=y(x_n)-y_n\) 满足

\[
\begin{cases}
\displaystyle l(e_n)=\frac{e_{n+1}-2e_n+e_{n-1}}{h^2}-q_ne_n=\frac{h^2}{12}y^{(4)}(\xi_n),\\[6pt]
e_0=e_N=0,
\end{cases}
\]

其中 \(\xi_n\) 一般是不知道的。于是转向下列差分问题：

\[
\begin{cases}
\displaystyle l(\varepsilon_n)=\frac{\varepsilon_{n+1}-2\varepsilon_n+\varepsilon_{n-1}}{h^2}-q_n\varepsilon_n=-\frac{h^2}{12}M,\\[6pt]
\varepsilon_0=\varepsilon_N=0,
\end{cases}
\tag{5.7.13}
\]

其中 \(M=\max_{a\le x\le b}|y^{(4)}(x)|\)。首先证明上述两个差分问题的解存在下列关系：

\[
|e_n|\le\varepsilon_n\qquad(n=0,1,\cdots,N).
\tag{5.7.14}
\]

事实上，由于

\[
l(\varepsilon_n)=-\frac{h^2}{12}M\le-\frac{h^2}{12}|y^{(4)}(\xi_n)|=-|l(e_n)|,
\]

故有

\[
l(\varepsilon_n-e_n)\le0,\qquad l(\varepsilon_n+e_n)\le0.
\]

又

\[
\varepsilon_0-e_0=\varepsilon_N-e_N=0,\qquad\varepsilon_0+e_0=\varepsilon_N+e_N=0,
\]

利用定理 5.2，知 \(\varepsilon_n-e_n\ge0,\varepsilon_n+e_n\ge0\)，即式（5.7.14）成立。

差分问题（5.7.13）的解仍不易直接求得，于是进一步考察

\[
\begin{cases}
\displaystyle\tilde l(\rho_n)=\frac{\rho_{n+1}-2\rho_n+\rho_{n-1}}{h^2}=-\frac{h^2}{12}M,\\[6pt]
\rho_0=\rho_N=0.
\end{cases}
\tag{5.7.15}
\]

这里 \(\tilde l(\rho_n-\varepsilon_n)=-q_n\varepsilon_n\le0\)，又 \(\rho_0-\varepsilon_0=\rho_N-\varepsilon_N=0\)，故由定理 5.2（注意，当 \(q_n=0\) 时，\(l(\tilde y_n)\) 就是 \(l(y_n)\)）知 \(\rho_n-\varepsilon_n\ge0\)，即 \(\varepsilon_n\le\rho_n\)，于是有 \(|e_n|\le\varepsilon_n\le\rho_n\)。

然而 \(\rho_n\) 是容易求出的。求解差分问题（5.7.15）所对应的边值问题

\[
\begin{cases}
\displaystyle\rho''=-\frac{h^2}{12}M,\\[6pt]
\rho(a)=\rho(b)=0,
\end{cases}
\]

解得

\[
\rho(x)=\frac{h^2}{24}M(x-a)(b-x).
\]

容易验证 \(\rho_n=\rho(x_n)\) 是差分问题（5.7.15）的解，注意到 \(\rho(x)\) 在点 \(x=\dfrac{a+b}2\) 达到最大值

\[
\rho\left(\frac{a+b}2\right)=\frac{M(b-a)^2}{96}h^2,
\]

因此有估计式（5.7.11）。证毕。

据估计式（5.7.11）知，当 \(h\to0\) 时有 \(e_n\to0\)，这表明差分方法（5.7.10）是收敛的。

\begin{quote}
校注（agent 补充，ch05b-E006）：式（5.7.15）后括注的原图在 \(y\) 上方印波浪号，即 \(l(\tilde y_n)\)，这里照录。按该式定义与``当 \(q_n=0\) 时''的上下文，此处疑将波浪号排错位置，预期为 \(\tilde l(y_n)\)，即去掉零阶项后的算子；原文并未定义 \(\tilde y_n\)。
\end{quote}

\subsection{本节覆盖原页的校注记录}\label{ux672cux8282ux8986ux76d6ux539fux9875ux7684ux6821ux6ce8ux8bb0ux5f55}

下列记录按原页关联，含同页相邻小节的校注；计算时同时读取上文校注和完整原页。

\begin{itemize}
\tightlist
\item
  \texttt{ch05b-E006}：原书 PDF 页 152
\end{itemize}

\end{document}
